\section{The Acting Person}

Before he became Pope John Paul II, \textbf{Karol Woyjtyla} wrote a phenomenological study of the \emph{acting person}. Although academic phenomenology is idealist, Woyjtyla employed its methodology in the service of gaining a deeper understanding of some aspects of Thomist metaphysical realism. In principle, it is the proper way to proceed; metaphysical speculation needs to be grounded in experience. A written text is the final stage of a Hermetic process that \textbf{Valentin Tomberg} describes:

\begin{enumerate}
\item The state of concentration without effort (beyond words and thoughts) 
\item Vigilant inner silence 
\item The inspired activity of imagination and thought 
\item The conscious self halts its creative activity and summarizes everything that preceded 
\end{enumerate}
The first text is mystical and symbolic. Then, it can be summarized as a system of metaphysics. We can see that in the development of the Vedanta. Rishis, or seers, at some point wrote down their experiences in the Vedas as poetry, myth, rites, laws, and so on. This was later summarized in the six philosophical schools which are based on the Vedas. But the metaphysical systems are not gnosis, rather they point the way back for us to reach that original revelation. Thus, a Hermetic phenomenology is deeper than academic phenomenology. Furthermore, it demands a change in the consciousness of the Hermetist, unlike a scientist who can learn his subject matter apart from any moral or spiritual development.

Woyjtyla's effort is promising although it falls short of leading us to a true gnosis. He focuses on the phenomenology of the person as an active agent and its relationship to the body. That will lead him to the spirit or intellectual soul\footnote{\url{https://gornahoor.net/?p=3390}}. It is important not to jump to a conclusion. Since metaphysics is the final stage in the process, it necessarily describes an accomplishment. Thus, Aquinas offers a proof of the hylomorphic unity of the human being. However, a careful observation of our inner states reveals something else. For example, we may notice inner conflicts, incompatible desires, or uncontrollable urges, for example, which belie any sense of unity. At this stage of our development, our unity is still virtual and not actual.

Thus, a fully developed phenomenology will not assume that unity at the outset, but will instead observe the activities of the intellectual, sensitive, and nutritive souls in isolation and in relationship to each other. A vigilant inner silence and concentration without effort are required to pick out these details in the stream of consciousness. This can be fruitful in many respects. One of them, which is scarcely dwelled upon, is that we can learn the states of consciousness of animals and plants. There are many stories\footnote{\url{http://www.anaflora.com/articles/saints-sages/index.html}} of the rapport that saints and sages have had with animals.

\paragraph{The Soul as Principle of Transcendence}
Wojtyla claims to have analyzed the person on the somatic and psychic level, yet, due to his refusal to distinguish the different aspects of the soul, he proceeds to an inconsistent conclusion regarding transcendence. He writes:

\begin{quotex}
Man had no direct experience of his soul. Experience of the transcendence of the person in the action together with all the elements and aspects of this experience is in no way equivalent to a direct experience of the soul… Both the reality itself of the soul and that of the soul's relation to the body are in this sense transphenomenal and extraexperiential. 

\end{quotex}
This is man as spirit, or man as a rational animal with an intellectual soul. Spirit as subject can never be an object of consciousness; hence, it is never part of our phenomenal world nor of our conscious experience. Yet Wojtyla hesitates to go to the next step. This comment seems hopeful:

\begin{quotex}
The soul-body relation is intuitively given in the experience of man as a real being. In this respect the subordination of the system of integration of the human person to the transcendence of the person in the action is revelatory. 

\end{quotex}
Intuition is a direct knowledge of being, beyond any conceptual knowledge. This is how man knows himself, through intuition, not as an object in the world, still less as the conclusion of a logical argument. To know is to be. Thus man knows himself to the extent that he is a real being. And he becomes a real being by actualizing his possibilities, by becoming a real man. Man's task then is to integrate himself, that is, to make the spirit the centre of his being as well as the dominant element. This is perfect, except we find this strange lapse, where Wojtyla writes:

\begin{quotex}
It is to metaphysical analysis that we owe the knowledge of the human soul as the principle underlying the unity of the being and the life of a concrete person. We infer the existence of the soul and its spiritual nature from effects that demand a sufficient reason, that is to say, a commensurate cause. In this perspective it is evident that there can be no such thing as a direct experience of the soul.

\end{quotex}
Of course, those aspects of the soul which belong to the spirit or intellectual soul — obligation, responsibility, truthfulness, self-determination, and consciousness — are not part of direct experience. Nevertheless, people do have experience of the lower soul through their thoughts, feelings, and desires. However, what is disappointing is the claim that we only infer the existence of the soul. What happened to ``know thyself", or even the intuition of the soul which Wojtyla also admits? We don't and shouldn't live our life by inferring we have a transcendent and eternal soul. We agree that the intuition of the intellectual soul in the mass of mankind for whom life lives through them is virtual, assuming it exists at all. However, we prefer the other formulation of the intuition of man's being which comes through his efforts at self-integration.



\flrightit{Posted on 2011-11-21 by Cologero }
